\documentclass[a4,11pt]{aleph-notas}
% Se puede ver la documentación aquí:
% https://github.com/alephsub0/LaTeX_aleph-notas

% -- Paquetes adicionales
\usepackage{enumitem}
\usepackage{url}
\usepackage{array}
\usepackage{booktabs}
\usepackage{longtable}
\usepackage{aleph-comandos}
\hypersetup{
    urlcolor=blue,
    linkcolor=blue,
}


% -- Datos
\institucion{Facultad de Ciencias Exactas, Naturales y Ambientales}
\carrera{Catálogo STEM}
\asignatura{Matemática III}
\tema{Plan tema 07: Campos vectoriales y regla de la cadena}
\autor{Andrés Merino}
\fecha{Periodo 2026-1}

\logouno[0.16\textwidth]{Logos/logoPUCE_04_ac}
\definecolor{colortext}{HTML}{1957d1}
\definecolor{colordef}{HTML}{1957d1}
\fuente{montserrat}


\begin{document}

\encabezado

%%%%%%%%%%%%%%%%%%%%%%%%%%%%%%%%%%%%%%%%
\section*{Resultado de Aprendizaje}
%%%%%%%%%%%%%%%%%%%%%%%%%%%%%%%%%%%%%%%%

%%%%%%%%%%%%%%%%%%%%%%%%%%%%%%%%%%%%%%%%
\subsection*{RdA de la asignatura:}
\begin{itemize}[leftmargin=*]
    \item \textbf{RdA 1:} Identificar conceptos, teoremas y propiedades de funciones vectoriales, derivadas parciales e integral vectorial.
    \item \textbf{RdA 2:} Calcular derivadas parciales e integrales con asistentes computacionales.
\end{itemize}

%%%%%%%%%%%%%%%%%%%%%%%%%%%%%%%%%%%%%%%%
\subsection*{RdA del tema:}
\begin{itemize}[leftmargin=*]
    \item Calcular la matriz jacobiana de un campo vectorial en un punto.
    \item Determinar si un campo vectorial es diferenciable en un punto.
    \item Aplicar la regla de la cadena para calcular derivadas parciales de composiciones de campos vectoriales.
    \item Relacionar la jacobiana de una composición con el producto de las jacobianas.
\end{itemize}

%%%%%%%%%%%%%%%%%%%%%%%%%%%%%%%%%%%%%%%%
\section*{Introducción}
%%%%%%%%%%%%%%%%%%%%%%%%%%%%%%%%%%%%%%%%

\paragraph{Pregunta inicial:}
Si un robot transforma coordenadas de posición y luego una cámara transforma esas posiciones en píxeles, ¿cómo cambia un píxel al mover ligeramente al robot? ¿Podría calcularse eso sin conocer directamente la función total?

%%%%%%%%%%%%%%%%%%%%%%%%%%%%%%%%%%%%%%%%
\section*{Desarrollo}
%%%%%%%%%%%%%%%%%%%%%%%%%%%%%%%%%%%%%%%%

%%%%%%%%%%%%%%%%%%%%%%%%%%%%%%%%%%%%%%%%
\subsection*{Actividad 1: Diferenciabilidad y matriz jacobiana}
Se extiende la diferenciabilidad a campos vectoriales y se introduce la matriz jacobiana como representante del diferencial.

\paragraph{¿Cómo lo haremos?}
\begin{itemize}[leftmargin=*]
    \item \textbf{Motivación:} retomar la pregunta inicial; mostrar que para funciones $\func{F}{\R^n}{\R^m}$ con $m>1$ la derivada ya no es un número ni un vector, sino una aplicación lineal representada por una matriz.
    \item \textbf{Clase magistral:} se definen derivadas parciales y jacobiana de un campo vectorial; se enuncia la diferenciabilidad y su relación con la jacobiana; se contrasta con el caso escalar. Se usa el resumen \href{https://andres-merino.github.io/MatematicaIII-05-N0051/01Resumen/Resumen07.pdf}{Resumen07.pdf}.
    \item \textbf{Implementación computacional:} calcular matrices jacobianas de distintos campos vectoriales con un asistente computacional.
    \item \textbf{Resolución de ejercicios:} calcular jacobianas y verificar diferenciabilidad, usando \href{https://andres-merino.github.io/MatematicaIII-05-N0051/04Ejercicios/Ejercicios07.pdf}{Ejercicios07.pdf}.
\end{itemize}

\paragraph{Verificación de aprendizaje:}
\begin{itemize}[leftmargin=*]
    \item ¿Cuál es la matriz jacobiana de $F(x,y)=(x^2+y,\, xy)$ en el punto $(1,2)$?
    \item ¿Por qué la existencia de la jacobiana no garantiza la diferenciabilidad de un campo vectorial?
    \item ¿Qué condición sobre las derivadas parciales garantiza que $F$ es diferenciable?
\end{itemize}

%%%%%%%%%%%%%%%%%%%%%%%%%%%%%%%%%%%%%%%%
\subsection*{Actividad 2: Regla de la cadena}
Se enuncia y aplica la regla de la cadena para composiciones de campos vectoriales diferenciables.

\paragraph{¿Cómo lo haremos?}
\begin{itemize}[leftmargin=*]
    \item \textbf{Clase magistral:} se enuncia la regla de la cadena $(F\circ G)'(a)=F'(G(a))\circ G'(a)$ y su versión en jacobianas $J_{F\circ G}(a)=J_F(G(a))\,J_G(a)$; se trabaja la notación de derivadas parciales compuestas. Se usa el resumen \href{https://andres-merino.github.io/MatematicaIII-05-N0051/01Resumen/Resumen07.pdf}{Resumen07.pdf}.
    \item \textbf{Resolución de ejercicios:} calcular derivadas parciales de composiciones usando la regla de la cadena, usando \href{https://andres-merino.github.io/MatematicaIII-05-N0051/04Ejercicios/Ejercicios07.pdf}{Ejercicios07.pdf}.
\end{itemize}

\paragraph{Verificación de aprendizaje:}
\begin{itemize}[leftmargin=*]
    \item ¿Qué relación hay entre la jacobiana de $F\circ G$ y las jacobianas de $F$ y $G$?
    \item Dados $G(s,t)=(s^2,st)$ y $F(x,y)=x+y$, ¿cuál es $\parcial{(F\circ G)}{s}$ en $(1,1)$?
    \item ¿En qué se diferencia la regla de la cadena para campos vectoriales de la versión escalar?
\end{itemize}

%%%%%%%%%%%%%%%%%%%%%%%%%%%%%%%%%%%%%%%%
\section*{Cierre}
%%%%%%%%%%%%%%%%%%%%%%%%%%%%%%%%%%%%%%%%

\paragraph{Tarea:}
Resolver del libro \href{https://catalogobiblioteca.puce.edu.ec/cgi-bin/koha/opac-detail.pl?biblionumber=83405&query_desc=su%3A%22An%C3%A1lisis%20vectorial%22}{Cálculo Vectorial de Colley}: sección 2.5, ejercicios 1--27 (impares).

\paragraph{Pregunta de investigación:}
\begin{enumerate}[leftmargin=*]
    \item ¿Cómo se usa la jacobiana en visión por computadora o en robótica para calcular transformaciones de movimiento?
    \item ¿Qué ocurre cuando la jacobiana de una transformación es singular (determinante cero)? ¿Qué implica geométricamente?
\end{enumerate}

\paragraph{Para el próximo tema:}
Ver los videos:
\begin{itemize}[leftmargin=*]
    \item \href{https://youtu.be/qb40J4N1fa4?si=edqBdY06q7YXj68v}{Diferenciación implícita, ¿qué está pasando aquí?}.
    \item \href{https://youtu.be/SUfJPLNZwBo?si=xHjCz3ZnXwdLOpaz}{Implicit Function Theorem}.
\end{itemize}


\end{document}
